\documentclass[aspectratio=169]{beamer}

\usepackage{graphicx} % Required for inserting images
\usepackage[english, bidi=basic, provide=*]{babel}

% Set default font to Sans Serif

% --- BEAMER THEME ---
\usetheme{Madrid}
\usecolortheme{default}
\setbeamertemplate{navigation symbols}{}

\title[VR Experiment EDA]{Initial EDA Report}
\subtitle{Headtracking Measures as Indicators of Depressive Symptoms \\ \small 360$^\circ$ VR Replication Study}
\author{Ankita Porel, Piyush Kumar, Shriram R}
\date{March 2026}

\begin{document}

\begin{frame}
    \titlepage
\end{frame}

\begin{frame}{Dataset Overview \& Demographics}
    \textbf{Inclusion \& Exclusion}
    \begin{itemize}
        \item Initial Sample: $N = 40$ college students.
        \item Exclusion: 3 participants removed due to significant simulator sickness (VRISE $< 25$).
        \item Final EDA Dataset: \textbf{$N = 37$}.
    \end{itemize}

    \vspace{0.5cm}
    \textbf{Demographics (N=37)}
    \begin{itemize}
        \item \textbf{Age:} $M = 22.89$ ($SD = 1.82$), Range: 19--27.
        \item \textbf{Gender:} 33 Male / 4 Female.
        \item \textbf{VR Experience:} Generally low ($M = 1.41$ on a 2-point scale).
    \end{itemize}
\end{frame}

\begin{frame}{Mental Health Score Distributions}
    \begin{columns}
        \column{0.4\textwidth}
        \begin{itemize}
            \item \textbf{PHQ-9:} Right-skewed ($M=5.81$). Sample is mostly in "Minimal" to "Mild" ranges.
            \item \textbf{GAD-7:} Matches depression skew ($M=5.11$).
            \item \textbf{STAI-T:} Shows wider variance ($SD=14.6$).
        \end{itemize}

        \column{0.6\textwidth}
        \begin{figure}[htbp]
          \centering
          \includegraphics[width=0.9\textwidth]{eda_plot1_distributions.png}
        \end{figure}
    \end{columns}
\end{frame}

\begin{frame}{Variable Covariance Analysis}
    \begin{columns}
        \column{0.4\textwidth}
        \textbf{Key Correlations:}
        \begin{itemize}
            \item \textbf{PHQ-9 \& GAD-7:} Strong positive correlation ($r=0.612$).
            \item \textbf{PHQ-9 \& STAI-T:} Strong positive correlation ($r=0.617$).
        \end{itemize}
        \vspace{0.5cm}
        \textit{Takeaway:} Anxiety must be treated as a covariate in future inferential models to isolate depressive psychomotor effects.

        \column{0.6\textwidth}
        \begin{figure}[htbp]
          \centering
          \includegraphics[width=0.9\textwidth]{eda_plot2_correlations.png}
        \end{figure}
    \end{columns}
\end{frame}

\begin{frame}{Subjective Affective Responses}
    \begin{columns}
        \column{0.4\textwidth}
        \textbf{Video Subjective Ratings}
        \begin{itemize}
            \item \textbf{Valence:} V2 (Beach) and V5 (Tahiti) rated highly pleasant. V4 (Nun) rated highly unpleasant.
            \item \textbf{Arousal:} V4 elicited the highest arousal spikes.
            \item Emotional manipulation via 360$^\circ$ video environments was successful.
        \end{itemize}

        \column{0.6\textwidth}
        \begin{figure}[htbp]
          \centering
          \includegraphics[width=0.9\textwidth]{eda_plot3_video_subjective.png}
        \end{figure}
    \end{columns}
\end{frame}

\begin{frame}{Objective Psychomotor Behavior by Video}
    \begin{columns}
        \column{0.4\textwidth}
        \textbf{Behavioral Shifts}
        \begin{itemize}
            \item \textbf{V1 (Abandoned):} Highest active exploration (Speed $M=39.6$).
            \item \textbf{V4 (Nun Horror):} Induced a "freezing" effect; lowest speed ($24.8$) and highly restricted Yaw scanning ($SD=33.9$).
            \item \textbf{V2/V5 (Water):} High horizontal scanning, low vertical variance.
        \end{itemize}

        \column{0.6\textwidth}
        \begin{figure}[htbp]
          \centering
          \includegraphics[width=0.9\textwidth]{eda_plot4_headtracking_by_video.png}
        \end{figure}
    \end{columns}
\end{frame}

\begin{frame}{Controlling for False Discovery Rate (FDR) -- Benjamini-Hochberg}
    \begin{itemize}
        \item Instead of controlling the probability of making \textit{even one} false positive (FWER), FDR controls the \textbf{expected proportion of false discoveries} among all significant results.
        \item Ideal for \textbf{exploratory science} (like this VR study), as it preserves statistical power while still keeping error rates in check.
    \end{itemize}

    \vspace{0.4cm}
    \textbf{The Benjamini-Hochberg (B-H) Procedure:}
    \begin{enumerate}
        \item \textbf{Rank} all $m$ obtained $p$-values from smallest to largest: $P_{(1)} \dots P_{(m)}$
        \item \textbf{Calculate} the B-H critical value for each test: $\frac{i}{m} \times Q$ \\
        \item \textbf{Evaluate:} Find the largest $p$-value that is smaller than its corresponding critical value. That test, and \textbf{all tests ranked above it}, are deemed statistically significant.
    \end{enumerate}

    \vspace{0.4cm}
    \begin{alertblock}{Why not Bonferroni?}
        Bonferroni ($\alpha/m$) is too stringent and assumes tests are independent. B-H provides a more nuanced, powerful threshold for highly correlated data (e.g., repeated headtracking measures).
    \end{alertblock}
\end{frame}


\begin{frame}{Conclusion \& Next Steps}
    \textbf{Summary of EDA}
    \begin{enumerate}
        \item Sample is strictly non-clinical but shows expected mental health score distributions.
        \item Headtracking metrics (Speed, Pitch, Yaw) are highly sensitive to the VR environment presented.
    \end{enumerate}

    \vspace{0.5cm}
    \textbf{Next Phases of Analysis}
    \begin{itemize}
        \item \textbf{Repeated Measures ANOVA} to test for significant differences in headtracking across the different video environments.
        \item Apply the \textbf{Benjamini-Hochberg (B-H) correction} to all resulting $p$-values to safely control the False Discovery Rate across our multiple comparisons.
    \end{itemize}
\end{frame}

\end{document}
